
\subsubsection{6.1 Interpretation of Results}

The experiments provide evidence that loss trajectory features can identify minority samples on the Waterbirds benchmark, and that this signal is associated with spurious correlation dynamics rather than generic sample difficulty. The finding that initial loss alone ($L_{1})$ achieves comparable performance to multi-feature trajectory analysis suggests that the discriminative information is present from the first epoch, rather than emerging through extended training dynamics.

The mechanism underlying this finding may be understood as follows: pretrained models encode visual features that are spuriously correlated with labels in the Waterbirds dataset (background features). When minority samples---where these features provide incorrect signal---enter training, they experience higher initial loss. GroupDRO reduces spurious reliance by upweighting minority groups, thereby reducing the initial loss gap; random reweighting only smooths gradients without affecting the underlying feature reliance.

The negative result on curvature timing refines understanding of the mechanism. The original hypothesis proposed that minority samples would show delayed curvature stabilization due to prolonged optimization conflict. The data do not support this: curvature stabilizes early (epoch 2-3) for all samples. The discriminative signal is magnitude-based (high $L_1$ for minority samples) rather than timing-based.

\subsubsection{6.2 Limitations}

Several limitations scope the claims that can be made from these experiments:

\textbf{Single dataset.} All experiments were conducted on Waterbirds. While this is a standard benchmark, generalization to other spurious correlation settings (CelebA, ColoredMNIST, real-world datasets) has not been established.

\textbf{Pretrained models only.} The finding that $L_1$ dominates may depend on pretrained features already encoding spurious patterns. Models trained from scratch may exhibit different dynamics.

\textbf{Detection, not intervention.} These experiments demonstrate that trajectory features can identify minority samples, not that this identification translates to successful debiasing. Prior work has shown that detection and intervention are distinct challenges.

\textbf{Limited seed variance for H-E1.} The existence test used a single seed with cross-validation. While H-M2 used three seeds and confirmed consistency, broader seed analysis would strengthen robustness claims.

\textbf{Observational design for specificity.} The specificity test uses GroupDRO as an intervention, but causal claims are limited. The comparison with variance-matched random reweighting provides stronger evidence than pure observation, but does not establish causation.

\textbf{Incomplete hypothesis coverage.} A planned experiment testing whether early trajectory divergence predicts final worst-group accuracy (H-M3) was not conducted, as it depended on the curvature timing mechanism that was not supported.

\subsubsection{6.3 Relation to Prior Work}

The results are consistent with prior findings that per-sample training dynamics carry meaningful information (Toneva et al., 2018; Li et al., 2025). The contribution here is demonstrating that such dynamics can specifically identify spurious correlation-affected samples, not just generally difficult samples. The controlled comparison with random reweighting provides evidence for this specificity that is absent from prior trajectory analysis work.

The finding that initial loss dominates is consistent with the observation that pretrained models already encode spurious features. This differs from the forgetting-based analysis of Toneva et al. (2018), where temporal patterns (repeated forgetting) were informative. For spurious correlation detection with pretrained models, the initial state appears more informative than subsequent dynamics.
